\section{Conclusion}\label{sec:conclusion}

This dissertation designed and audited a
\emph{partial-observation routing-influence filter}. The filter ingests a
road graph, sparse regime-labelled station observations at $\Xo$, and
held-out trajectories; it fits one inverse-Markov-chain model per regime;
then it scores each trajectory with a log-likelihood ratio. The underlying
estimator is Morimura et al.'s inverse-MC method~\cite{morimura2013}; the
contribution here is the two-regime detector built around it, the
trajectory-level scoring rule, and the diagnostic suite used to decide
whether a positive AUC is actually evidence of route-choice recovery.

\paragraph{What the evidence supports.}
Tier~1 validates the estimator implementation: on Morimura et al.'s
synthetic recovery benchmark, the reproduced RMAE curve tracks the
published proposed-method curve for $|\Xo|\ge10$, with analytic gradients
finite-difference checked before downstream use. Tier~2 is the positive
detector result. In labelled SUMO simulation, with shared OD pairs and a
known rerouting flag, the filter reaches $\AUC=0.708$ at $T=40$ against a
full-observation empirical $1$-step reference of $0.801$. The
feature-regularised variant does not robustly lift aggregate AUC across
all seeds, but it does robustly improve per-row chain recovery, which is
the stronger claim. Tier~3 is the real-world stress test. On Xuancheng AVI
data, where only behavioural proxy labels are available, the fitted
$1$-step detector remains near chance after OD control: rush-vs-off-peak
does not recover at any tested OD granularity, and the stronger
Labour-Day-vs-normal split falls to $\AUC=0.502$ under
validation-selected multistart despite a weak full-observation empirical
signal.

\paragraph{What the audits changed.}
Several headline-looking results became more precise once audited. The
SUMO features result became a chain-recovery result rather than a robust
detection-AUC result. The SUMO $2$-step empirical backoff lifted AUC by
about $+0.09$--$+0.10$ across all five seeds, showing that higher-order
route memory is a real signal in the controlled data and a justified
future model class. The same $2$-step audit did not help on the Xuancheng
Labour-Day split, so route memory is not the simple missing lever for that
real-world proxy. The optional $f+g$ route is mathematically verified and
useful on small subgraphs when $\gamma$ is tuned near the $f$-only end, but
naive city-scale sparse-LU is too slow and empirical hitting rates are too
floor-dominated to use as the deployed full-network detector in this
thesis.

\paragraph{Where the filter is reliable.}
The filter is reliable as a controlled labelled detector under the
conditions tested in SUMO: matched OD pairs, known rerouting labels, and a
finite-data empirical reference against which chain recovery can be
audited. In that setting it recovers a substantial fraction of the
available $1$-step signal and exposes where the $1$-step model class runs
out of capacity.

\paragraph{Where it is not reliable.}
The filter does not support a positive real-world sat-navigation detection
claim on Xuancheng. Without trip-level guidance/adoption labels,
rush-vs-off-peak and Labour-Day splits remain behavioural proxies that mix
route guidance, ordinary congestion adaptation, driver composition, OD
mix, time heterogeneity and data-quality effects. Once those confounds are
audited, the fitted partial-observation detector does not robustly recover
the within-OD signal that remains in the full-observation reference.

\begin{figure}[H]
\centering
\includegraphics[width=0.92\linewidth]{fig08_conclusion_synthesis.png}
\caption{Cross-tier synthesis of the dissertation's headline results.
The bars use a normalised summary, since the tiers use different metrics:
Tier~1 reports benchmark-reproduction closeness, computed as the mean
capped ratio between the visual paper RMAE and this reproduction's RMAE for
$|\Xo|\ge10$; Tiers~2--3 report AUC gap closed,
$(\AUC_{\mathrm{fit}}-0.5)/(\AUC_{\mathrm{emp}}-0.5)$. The central pattern
is that the estimator reproduces its synthetic benchmark, the labelled
SUMO filter recovers a substantial fraction of the empirical $1$-step
signal, and the real Xuancheng behavioural proxy does not recover under
the fitted partial-observation model.}
\label{fig:conclusion-synthesis}
\end{figure}

The practical conclusion is therefore sharper than a single AUC number:
the inverse-MC machinery is not the main bottleneck for real-world
externally-influenced routing detection. The bottleneck is validation
data. A future deployment needs at least partial trip-level
sat-navigation-use or guidance-compliance labels, repeated station and OD
sensitivity checks, and a richer model class where the controlled data
show it is warranted. Until those labels exist, this framework is best
read as a validated controlled-simulation filter and a real-world
confounding audit, not as a finished real-world sat-navigation detector.
